\section{安全目标建模：从 CIA 到 Agent 行为约束}

课程把传统 CIA 目标引入 Agent 场景，并强调它们需要被 ``行为化''。例如 confidentiality 不只是防止数据库泄露，还包括防止 Agent 通过工具链间接泄露敏感上下文；integrity 不只是防篡改，还包括保证推理链和动作链不被污染；availability 不只是服务在线，还包括防止攻击者用高成本任务拖垮执行资源。

\begin{importantbox}{Agent 时代的 CIA 扩展}
\begin{itemize}
    \item \textbf{Confidentiality}：敏感数据不得被未授权读取、推断或外发。
    \item \textbf{Integrity}：状态、计划、工具调用与执行结果不可被隐蔽篡改。
    \item \textbf{Availability}：系统在对抗流量、恶意任务和资源争用下保持可服务。
    \item \textbf{Action Safety}：即使被诱导，Agent 也不能执行高危未授权动作。
\end{itemize}
\end{importantbox}

讲者还比较了传统系统和 Agent 系统在安全目标上的区别。传统系统通常功能边界清晰、权限边界固定；Agent 系统在运行中会持续扩展上下文和操作空间，因此目标约束必须是动态的，不能只靠静态 ACL。

\begin{warningbox}{静态策略的局限}
如果策略仅按 ``系统角色'' 固定授权，而不结合 ``当前任务上下文''，Agent 很容易在一次合法会话中被间接注入后执行不该执行的动作。
\end{warningbox}

\subsection{本章小结}
安全目标需要从 ``对象安全'' 升级为 ``行为安全''。在 Agent 场景下，动态上下文策略和执行时约束比静态访问控制更关键。

